% Exploratory addendum. Tables are generated from saved results (addendum/v1/report,
% addendum/v2/report). Prose drafted with AI assistance for author editing; every number
% traces to addendum/v1/report/summary.json or addendum/v2/report/summary.json.
\subsection{Exploratory addendum}
\label{sec:addendum}

The experiments in this subsection were declared after the main study, in two settings files committed before any score was computed, and each was executed once. They test four gaps identified in review: the small number of datasets, the half-budget minority quota, the absence of an equal-budget B2 comparison and the absence of published curation and oversampling baselines. None of these results enters the preceding tables, and all tests are exploratory, without multiplicity adjustment or significance claims.

\begin{table}[!t]
\centering
\scriptsize
\setlength{\tabcolsep}{3pt}
\caption{Exploratory addendum, declared after the main study. Datasets. IR = majority/minority; $b$ = selection budget; TEST positives are per split (five stratified 80/20 splits).}
\label{tab:addendum-datasets}
\resizebox{\linewidth}{!}{%
\begin{tabular}{lrrrcrcl}
\toprule
Dataset & Rows & Min. & IR & Nom. & $b$ & TEST pos. & Pool gen. \\
\midrule
Abalone19 & 4174 & 32 & 129.4 & yes & 56 & 6 & SMOTENC \\
Ecoli4 & 336 & 20 & 15.8 & no & 70 & 4 & SMOTE \\
page-blocks0 & 5472 & 559 & 8.8 & no & 49 & 112 & RandomOverSampler \\
yeast3 & 1484 & 163 & 8.1 & no & 65 & 33 & SMOTE \\
yeast4 & 1484 & 51 & 28.1 & no & 30 & 10 & SMOTE \\
vehicle0 & 846 & 199 & 3.3 & no & 70 & 40 & SMOTE \\
segment0 & 2308 & 329 & 6.0 & no & 70 & 66 & SMOTE \\
ecoli1 & 336 & 77 & 3.4 & no & 31 & 15 & SMOTE \\
abalone9-18 & 731 & 42 & 16.4 & yes & 30 & 8 & SMOTENC \\
winequality-red-4 & \multicolumn{7}{l}{unavailable (empty source archive)} \\
\bottomrule
\end{tabular}}
\end{table}

Six further KEEL datasets join the three original ones (Table~\ref{tab:addendum-datasets}). They follow the fresh-split design of RQ6: five stratified splits, train-only fitting, minority-only SMOTE or SMOTENC pools of size $4b$, and the budget $b=\min(70,\max(30,\mathrm{round}(0.5\,n_{\mathrm{min}})))$. A seventh candidate, winequality-red-4, is unavailable because its source archive is empty. vehicle0 declares all 18 features as integers; unrounded SMOTE interpolation violates that declaration in all 1{,}400 candidates across its five splits, and bounded repair recovers only six, because nearly every record needs more than three integer edits. No curator can therefore run on vehicle0, although training comparators still can. This is the first fresh-split case in this study where the structural checks change the outcome: interpolation oversamplers can break source declarations wholesale, and the deliberately narrow repair bound declines to rewrite such records.

\begin{table}[!t]
\centering
\scriptsize
\setlength{\tabcolsep}{3pt}
\caption{Exploratory addendum, declared after the main study. Addendum v1 across-dataset summary: dataset-level mean difference and split wins/ties/losses over five fresh splits.}
\label{tab:addendum-v1-e4}
\resizebox{\linewidth}{!}{%
\begin{tabular}{lllrrr}
\toprule
Comparison & Clf & Metric & Datasets $+/-/0$ & Splits W/T/L & Median diff. \\
\midrule
QC $-$ adm. random & LR & AP & 3/4/1 & 17/7/16 & -0.001 \\
QC $-$ adm. random & LR & F1 & 4/1/3 & 11/22/7 & +0.001 \\
QC $-$ adm. random & RF & AP & 5/3/0 & 19/6/15 & +0.002 \\
QC $-$ adm. random & RF & F1 & 6/2/0 & 15/15/10 & +0.013 \\
QC $-$ real only & LR & AP & 3/5/0 & 18/6/16 & -0.000 \\
QC $-$ real only & LR & F1 & 5/1/2 & 15/15/10 & +0.006 \\
QC $-$ real only & RF & AP & 5/3/0 & 17/5/18 & +0.002 \\
QC $-$ real only & RF & F1 & 5/2/1 & 18/11/11 & +0.027 \\
Best aug. $-$ real only$^\ast$ & LR & AP & 6/2/1 & 25/9/11 & +0.003 \\
Best aug. $-$ real only$^\ast$ & LR & F1 & 6/1/2 & 18/17/10 & +0.006 \\
Best aug. $-$ real only$^\ast$ & RF & AP & 8/1/0 & 23/6/16 & +0.021 \\
Best aug. $-$ real only$^\ast$ & RF & F1 & 7/2/0 & 28/9/8 & +0.035 \\
\bottomrule
\end{tabular}}

\vspace{2pt}\parbox{\linewidth}{\scriptsize $^\ast$Best method chosen on the same TEST splits; optimistic by construction.}
\end{table}

Across the eight datasets with usable pools, quality-coverage does not outperform random selection from the same admitted pool (Table~\ref{tab:addendum-v1-e4}). The median dataset-level F1 difference is $+.001$ for logistic regression and $+.013$ for random forest; over 40 paired splits per classifier, wins, ties and losses are 11/22/7 and 15/15/10, and exploratory Wilcoxon $p$-values for F1 and AP range from $.19$ to $1.00$. The geometric effect, in contrast, replicates: on all five new datasets with usable pools, quality-coverage covers more distinct TRAIN support at radius $.05$ and selects fewer near-duplicates than admitted random selection, for example $.677$ versus $.461$ coverage and $.000$ versus $.327$ near-duplication on abalone9-18. Choosing the best augmentation per dataset can appear to beat real-only training, but that choice uses the same TEST splits and is optimistic by construction.

\begin{table*}[!t]
\centering
\scriptsize
\setlength{\tabcolsep}{3pt}
\caption{Exploratory addendum, declared after the main study. Addendum v1 historical-split designs at fixed budgets: canonical quality-coverage (QC, quota $0.5b$), QC with the real-ratio quota (E2), admitted random (mean of five seeds) and random subsets of the B2-accepted pool (E3, mean of five seeds).}
\label{tab:addendum-v1-e2e3}
\resizebox{\textwidth}{!}{%
\begin{tabular}{llrrrrrrrr}
\toprule
Dataset & Design & Quota & Cov. & Near dup. & LR TSTR F1 & LR Hyb. F1 & RF TSTR F1 & RF Hyb. F1 & LR Hyb. AP \\
\midrule
Abalone19 & QC & 28 & .948 & .000 & .031 & .000 & .030 & .000 & .017 \\
Abalone19 & QC real ratio (E2) & 0 & .948 & .000 & .029 & .000 & .027 & .000 & .018 \\
Abalone19 & Adm. random & -- & .833 & .493 & .022 & .000 & .030 & .000 & .021 \\
Abalone19 & B2 random (E3) & -- & .853 & .232 & .026 & .000 & .030 & .000 & .015 \\
\midrule
Ecoli4 & QC & 35 & .506 & .000 & .205 & 1.000 & .296 & 1.000 & 1.000 \\
Ecoli4 & QC real ratio (E2) & 4 & .506 & .000 & .190 & 1.000 & .250 & 1.000 & 1.000 \\
Ecoli4 & Adm. random & -- & .439 & .197 & .217 & .933 & .231 & 1.000 & 1.000 \\
Ecoli4 & B2 random (E3) & -- & .419 & .049 & .213 & .933 & .271 & 1.000 & 1.000 \\
\midrule
page-blocks0 & QC & 24 & .932 & .327 & .446 & .715 & .342 & .863 & .744 \\
page-blocks0 & QC real ratio (E2) & 5 & .932 & .000 & .593 & .722 & .631 & .873 & .747 \\
page-blocks0 & Adm. random & -- & .914 & .702 & .535 & .715 & .475 & .869 & .743 \\
page-blocks0 & B2 random (E3) & -- & .905 & .506 & .421 & .715 & .438 & .867 & .744 \\
\bottomrule
\end{tabular}}
\end{table*}

On the historical split (Table~\ref{tab:addendum-v1-e2e3}), setting the minority quota to $\lfloor b\times$ TRAIN minority fraction$\rfloor$ lowers it from 28, 35 and 24 to 0, 4 and 5 records. Coverage is unchanged on all three datasets. Pure TSTR F1 rises on page-blocks0 ($.446$ to $.593$ with logistic regression, $.342$ to $.631$ with random forest), where near-duplication also falls from $.327$ to zero, but it falls on Ecoli4 ($.296$ to $.250$ with random forest) and slightly on Abalone19. Hybrid F1 moves by at most $.011$. The half-budget quota therefore shapes synthetic-only results in both directions and is not the main determinant of hybrid utility. At an equal budget, quality-coverage covers more TRAIN support than random subsets of the B2-accepted, structurally complete pool on all three datasets ($.948$, $.506$ and $.932$ versus $.853$, $.419$ and $.905$), whereas F1 shows no consistent winner: B2 subsets are ahead on page-blocks0 random-forest TSTR ($.438$ versus $.342$) and behind on Ecoli4 ($.271$ versus $.296$).

\begin{table*}[!t]
\centering
\scriptsize
\setlength{\tabcolsep}{3pt}
\caption{Exploratory addendum, declared after the main study. Mean F1 and AP of the eight curators (hybrid training: real TRAIN plus $b$ selected candidates).}
\label{tab:addendum-curators}
\resizebox{\textwidth}{!}{%
\begin{tabular}{lrrrrrrrr}
\toprule
Dataset & Schema rand. & Adm. rand. & Source order & Farthest-first & Cov.-greedy & QC (ours) & CLLM-style$^\dagger$ & Disc.-realism$^\dagger$ \\
\midrule
\multicolumn{9}{l}{\textit{LR, mean F1 over five splits}} \\
Abalone19 & \textbf{.000} & \textbf{.000} & \textbf{.000} & \textbf{.000} & \textbf{.000} & \textbf{.000} & \textbf{.000} & \textbf{.000} \\
Ecoli4 & .758 & .758 & .758 & .758 & .758 & .758 & \textbf{.826} & .758 \\
page-blocks0 & .736 & \textbf{.741} & .738 & .738 & .740 & .738 & .736 & .739 \\
yeast3 & .773 & .769 & .779 & .773 & \textbf{.782} & \textbf{.782} & .771 & .779 \\
yeast4 & .152 & .150 & .154 & \textbf{.203} & .158 & .158 & .180 & .185 \\
segment0 & \textbf{.988} & \textbf{.988} & \textbf{.988} & \textbf{.988} & \textbf{.988} & \textbf{.988} & \textbf{.988} & \textbf{.988} \\
ecoli1 & .760 & .745 & .749 & .750 & .762 & .758 & .742 & \textbf{.762} \\
abalone9-18 & .676 & .677 & \textbf{.684} & .670 & .670 & .679 & .671 & .659 \\
\midrule
\multicolumn{9}{l}{\textit{LR, mean AP over five splits}} \\
Abalone19 & .041 & .047 & .046 & .047 & .045 & .044 & \textbf{.050} & .042 \\
Ecoli4 & \textbf{.944} & .934 & .934 & .934 & .934 & .934 & \textbf{.944} & .937 \\
page-blocks0 & .800 & .799 & .799 & .800 & .799 & .800 & \textbf{.800} & .799 \\
yeast3 & .813 & .809 & .807 & .809 & .807 & .808 & \textbf{.813} & .813 \\
yeast4 & .311 & .319 & .311 & .327 & .306 & .305 & \textbf{.331} & .322 \\
segment0 & .977 & .977 & .977 & .978 & .977 & \textbf{.978} & .974 & .977 \\
ecoli1 & .861 & .862 & \textbf{.868} & .864 & .867 & .865 & .864 & .861 \\
abalone9-18 & .774 & .780 & .771 & .780 & .781 & .775 & \textbf{.783} & .780 \\
\midrule
\multicolumn{9}{l}{\textit{RF, mean F1 over five splits}} \\
Abalone19 & .000 & \textbf{.057} & .000 & .000 & .000 & .000 & .000 & .000 \\
Ecoli4 & .786 & .748 & \textbf{.819} & .800 & .762 & .786 & .743 & .790 \\
page-blocks0 & .890 & \textbf{.893} & .892 & .891 & .886 & .885 & .890 & .885 \\
yeast3 & .775 & .790 & .787 & .794 & \textbf{.810} & .801 & .781 & .794 \\
yeast4 & .246 & .245 & .240 & .238 & .210 & .268 & \textbf{.335} & .199 \\
segment0 & .991 & .989 & .989 & \textbf{.991} & \textbf{.991} & \textbf{.991} & .989 & .989 \\
ecoli1 & .775 & .736 & .771 & \textbf{.782} & .773 & .754 & .748 & .762 \\
abalone9-18 & .348 & .373 & .402 & .370 & \textbf{.406} & .388 & .313 & .392 \\
\midrule
\multicolumn{9}{l}{\textit{RF, mean AP over five splits}} \\
Abalone19 & .040 & \textbf{.070} & .034 & .030 & .037 & .033 & .039 & .041 \\
Ecoli4 & .884 & .872 & \textbf{.895} & .882 & .870 & .882 & .811 & .883 \\
page-blocks0 & .922 & .922 & .924 & .925 & .924 & .926 & \textbf{.929} & .924 \\
yeast3 & .859 & .867 & .855 & .873 & .858 & .859 & .875 & \textbf{.875} \\
yeast4 & .406 & .379 & .401 & .415 & .364 & .407 & .383 & \textbf{.421} \\
segment0 & 1.000 & 1.000 & \textbf{1.000} & 1.000 & 1.000 & 1.000 & 1.000 & 1.000 \\
ecoli1 & .867 & .862 & .857 & .860 & .840 & .820 & .868 & \textbf{.870} \\
abalone9-18 & .480 & .533 & .511 & .551 & \textbf{.567} & .554 & .497 & .518 \\
\bottomrule
\end{tabular}}

\vspace{2pt}\parbox{\textwidth}{\scriptsize Bold: best in row. $^\dagger$Re-implementation of the published principle, not the authors' code. vehicle0 omitted: no structurally complete candidate.}
\end{table*}

The second addendum benchmarks quality-coverage inside the same pipeline, on the nine available datasets, the same splits, pools and classifiers. Two curators re-implement published principles rather than the authors' code: a learning-dynamics filter in the spirit of CLLM \cite{seedat2024cllm}, which keeps low-aleatoric, high-confidence candidates under a gradient-boosting model fitted on real TRAIN, and a discriminator-realism filter following the realism component of RUBRIC \cite{yu2026rubric} and sample-level auditing \cite{alaa2022faithful}, which keeps the candidates a random forest finds hardest to tell apart from real minority rows. All 45 rebuilt pools match the content fingerprints recorded earlier, and the 930 overlapping rows reproduce exactly. Among the eight curators (Table~\ref{tab:addendum-curators}), quality-coverage has the best average rank only for logistic-regression F1 ($3.75$ of 8); the CLLM-style filter ranks best for logistic-regression AP ($2.44$) and the discriminator for random-forest AP ($2.88$), and Friedman tests do not separate the curators ($p=.21$ to $.65$).

\begin{table*}[!t]
\centering
\scriptsize
\setlength{\tabcolsep}{3pt}
\caption{Exploratory addendum, declared after the main study. Mean F1 and AP of real-only training, class weighting, existing oversamplers, four published oversamplers (imbalanced-learn 0.14.2; full balancing and $+b$ rows) and quality-coverage.}
\label{tab:addendum-oversamplers}
\resizebox{\textwidth}{!}{%
\begin{tabular}{lrrrrrrrrrrrrr}
\toprule
Dataset & Real only & CW-LR & ROS+$b$ & SMOTE(NC)+$b$ & SMOTE-ENN full & SMOTE-ENN +$b$ & SMOTE-Tomek full & SMOTE-Tomek +$b$ & B-SMOTE full & B-SMOTE +$b$ & ADASYN full & ADASYN +$b$ & QC (ours) \\
\midrule
\multicolumn{14}{l}{\textit{LR, mean F1 over five splits}} \\
Abalone19 & .000 & \textbf{.051} & .000 & .000 & -- & -- & -- & -- & -- & -- & -- & -- & .000 \\
Ecoli4 & \textbf{.848} & .695 & .745 & .758 & .711 & .758 & .745 & .758 & .700 & .723 & .688 & .740 & .758 \\
page-blocks0 & .728 & .708 & .738 & .736 & .693 & \textbf{.761} & .699 & .743 & .665 & .738 & .628 & -- & .738 \\
yeast3 & .776 & .660 & .770 & .778 & .648 & .779 & .676 & .776 & .634 & \textbf{.782} & .641 & .776 & .782 \\
yeast4 & .101 & .235 & .147 & .154 & .209 & .099 & .231 & .200 & \textbf{.282} & .201 & .217 & .222 & .158 \\
vehicle0 & .927 & .923 & .931 & .929 & .892 & .907 & .922 & .931 & .919 & .930 & .918 & \textbf{.934} & -- \\
segment0 & .988 & .965 & .988 & .988 & \textbf{.991} & .988 & \textbf{.991} & .988 & .931 & .986 & .938 & .989 & .988 \\
ecoli1 & .751 & .746 & .754 & .749 & .749 & .737 & .749 & .745 & \textbf{.760} & .730 & .757 & .744 & .758 \\
abalone9-18 & .672 & .471 & .666 & .636 & -- & -- & -- & -- & -- & -- & -- & -- & \textbf{.679} \\
\midrule
\multicolumn{14}{l}{\textit{LR, mean AP over five splits}} \\
Abalone19 & .042 & \textbf{.045} & .040 & .041 & -- & -- & -- & -- & -- & -- & -- & -- & .044 \\
Ecoli4 & .944 & .923 & .908 & .944 & .896 & \textbf{.947} & .901 & .940 & .897 & .922 & .903 & .903 & .934 \\
page-blocks0 & \textbf{.800} & .788 & .800 & .800 & .789 & .794 & .790 & .800 & .752 & .800 & .748 & -- & .800 \\
yeast3 & .813 & .809 & .807 & .809 & .805 & .811 & .807 & \textbf{.814} & .795 & .810 & .801 & .808 & .808 \\
yeast4 & .313 & .251 & .308 & .309 & .242 & \textbf{.338} & .259 & .329 & .312 & .323 & .254 & .307 & .305 \\
vehicle0 & .982 & .976 & .980 & .981 & .927 & .921 & \textbf{.984} & .981 & .981 & .983 & .980 & .981 & -- \\
segment0 & .976 & .981 & .977 & .977 & .985 & .973 & .986 & .977 & \textbf{.997} & .983 & .996 & .983 & .978 \\
ecoli1 & .866 & .870 & \textbf{.874} & .870 & .842 & .842 & .867 & .867 & .870 & .872 & .868 & .865 & .865 \\
abalone9-18 & .771 & \textbf{.777} & .772 & .766 & -- & -- & -- & -- & -- & -- & -- & -- & .775 \\
\midrule
\multicolumn{14}{l}{\textit{RF, mean F1 over five splits}} \\
Abalone19 & \textbf{.000} & -- & \textbf{.000} & \textbf{.000} & -- & -- & -- & -- & -- & -- & -- & -- & \textbf{.000} \\
Ecoli4 & .600 & -- & .743 & .819 & \textbf{.836} & .776 & .800 & .786 & .724 & .743 & .798 & .798 & .786 \\
page-blocks0 & .893 & -- & .885 & \textbf{.901} & .830 & .856 & .865 & .895 & .864 & .896 & .850 & -- & .885 \\
yeast3 & .775 & -- & .783 & .777 & .778 & .771 & .804 & .786 & .792 & \textbf{.807} & .788 & .795 & .801 \\
yeast4 & .236 & -- & .243 & .165 & .329 & .103 & \textbf{.433} & .185 & .425 & .234 & .405 & .188 & .268 \\
vehicle0 & .928 & -- & .927 & .927 & .877 & .884 & .920 & .928 & .924 & \textbf{.931} & .907 & .928 & -- \\
segment0 & .992 & -- & .991 & .991 & .989 & .989 & .991 & .992 & .992 & .994 & .994 & \textbf{.994} & .991 \\
ecoli1 & .726 & -- & .771 & .738 & .787 & .772 & \textbf{.820} & .760 & .806 & .789 & .797 & .776 & .754 \\
abalone9-18 & .347 & -- & .373 & .374 & -- & -- & -- & -- & -- & -- & -- & -- & \textbf{.388} \\
\midrule
\multicolumn{14}{l}{\textit{RF, mean AP over five splits}} \\
Abalone19 & .029 & -- & .026 & .028 & -- & -- & -- & -- & -- & -- & -- & -- & \textbf{.033} \\
Ecoli4 & .828 & -- & .866 & .841 & \textbf{.925} & .846 & .880 & .878 & .884 & .859 & .910 & .863 & .882 \\
page-blocks0 & .925 & -- & .920 & .924 & .915 & .902 & .917 & .921 & .912 & .925 & .909 & -- & \textbf{.926} \\
yeast3 & .868 & -- & .851 & \textbf{.876} & .850 & .836 & .849 & .858 & .810 & .860 & .808 & .849 & .859 \\
yeast4 & .387 & -- & .364 & .349 & .352 & .332 & .421 & .403 & \textbf{.435} & .382 & .368 & .355 & .407 \\
vehicle0 & \textbf{.984} & -- & .982 & .983 & .931 & .951 & .974 & .982 & .980 & .980 & .975 & .979 & -- \\
segment0 & 1.000 & -- & 1.000 & 1.000 & 1.000 & .992 & 1.000 & 1.000 & 1.000 & 1.000 & 1.000 & \textbf{1.000} & 1.000 \\
ecoli1 & .833 & -- & .864 & .853 & .738 & \textbf{.875} & .848 & .855 & .800 & .829 & .846 & .844 & .820 \\
abalone9-18 & .545 & -- & \textbf{.557} & .522 & -- & -- & -- & -- & -- & -- & -- & -- & .554 \\
\bottomrule
\end{tabular}}

\vspace{2pt}\parbox{\textwidth}{\scriptsize Bold: best in row. --: unavailable (no nominal path, library error or no structurally complete candidate; CW-LR is logistic only). SMOTE-ENN/Tomek may remove rows; final counts are in results.}
\end{table*}

Four published oversamplers from imbalanced-learn \cite{lemaitre2017imbalanced} serve as training comparators: SMOTE-ENN and SMOTE-Tomek \cite{batista2004study}, Borderline-SMOTE \cite{han2005borderline} and ADASYN \cite{he2008adasyn}, each with full balancing and with $b$ added rows (Table~\ref{tab:addendum-oversamplers}). Full balancing helps random forest most on the hardest datasets, for example SMOTE-Tomek reaches F1 $.433$ on yeast4 against $.268$ for quality-coverage, but the same oversamplers often reduce logistic-regression F1, to $.63$--$.68$ on yeast3 against $.78$. Availability is itself a finding. Borderline-SMOTE and ADASYN offer no path for nominal features, the cleaning step of SMOTE-ENN and SMOTE-Tomek fails on nominal columns, and ADASYN with $b$ added rows generates no samples on page-blocks0; with vehicle0 lacking curators, the complete-case comparison covers five datasets. As context only, our real-only and SMOTE AUROC values exceed those reported by Newaz and Haq \cite{newaz2022hybrid} for six shared datasets under 10-fold cross-validation, and the sign of the SMOTE effect agrees with theirs in five of twelve dataset--classifier pairs.

\begin{table}[!t]
\centering
\scriptsize
\setlength{\tabcolsep}{3pt}
\caption{Exploratory addendum, declared after the main study. Average ranks (1 = best) of mean F1 and AP across datasets, with exploratory Friedman tests, and Wilcoxon signed-rank tests of quality-coverage against each method. No multiplicity adjustment or significance claim.}
\label{tab:addendum-ranks}
\resizebox{\linewidth}{!}{%
\begin{tabular}{lrrrr}
\toprule
\multicolumn{5}{l}{\textit{(a) Curators, 8 complete datasets}} \\
Method & LR F1 & LR AP & RF F1 & RF AP \\
\midrule
Schema rand. & 5.12 & 4.81 & 4.81 & 5.00 \\
Adm. rand. & 5.12 & 5.12 & 4.50 & 5.12 \\
Source order & 4.38 & 5.38 & 4.00 & 4.62 \\
Farthest-first & 4.69 & 3.50 & 3.62 & 4.25 \\
Cov.-greedy & 3.81 & 5.00 & 3.88 & 5.88 \\
QC (ours) & 3.75 & 4.88 & 4.06 & 4.25 \\
CLLM-style$^\dagger$ & 5.12 & 2.44 & 5.88 & 4.00 \\
Disc.-realism$^\dagger$ & 4.00 & 4.88 & 5.25 & 2.88 \\
Friedman $p$ & 0.654 & 0.207 & 0.525 & 0.387 \\
\midrule
\multicolumn{5}{l}{\textit{(b) All methods, 5 complete datasets}} \\
Method & LR F1 & LR AP & RF F1 & RF AP \\
\midrule
QC (ours) & 6.10 & 9.40 & 7.60 & 8.60 \\
Adm. rand. & 10.30 & 9.80 & 10.80 & 9.00 \\
Farthest-first & 7.30 & 7.60 & 6.60 & 6.80 \\
CLLM-style$^\dagger$ & 8.90 & 6.80 & 11.20 & 7.80 \\
Disc.-realism$^\dagger$ & 5.70 & 7.60 & 9.80 & 4.40 \\
Real only & 7.30 & 7.40 & 12.50 & 9.00 \\
ROS+$b$ & 9.80 & 9.20 & 10.00 & 8.60 \\
SMOTE(NC)+$b$ & 8.30 & 7.00 & 10.70 & 10.20 \\
SMOTE-ENN full & 8.50 & 13.20 & 7.20 & 10.40 \\
SMOTE-ENN +$b$ & 9.70 & 7.60 & 13.00 & 12.00 \\
SMOTE-Tomek full & 7.20 & 10.20 & 3.70 & 7.80 \\
SMOTE-Tomek +$b$ & 8.10 & 5.40 & 9.80 & 7.00 \\
B-SMOTE full & 10.00 & 9.00 & 6.00 & 7.60 \\
B-SMOTE +$b$ & 10.20 & 5.60 & 6.50 & 8.80 \\
ADASYN full & 10.80 & 9.90 & 4.50 & 8.20 \\
ADASYN +$b$ & 7.80 & 10.30 & 6.10 & 9.80 \\
Friedman $p$ & 0.904 & 0.578 & 0.031 & 0.826 \\
\midrule
\multicolumn{5}{l}{\textit{(c) Wilcoxon, QC minus method: datasets QC better/worse, $p$}} \\
Method & LR F1 & LR AP & RF F1 & RF AP \\
\midrule
Schema rand. & 4/1, 0.12 & 4/4, 0.64 & 4/2, 0.31 & 5/3, 0.95 \\
Adm. rand. & 4/1, 0.19 & 3/4, 0.58 & 6/2, 0.31 & 5/3, 1.00 \\
Source order & 4/1, 0.44 & 4/3, 0.81 & 3/4, 0.58 & 4/4, 0.84 \\
Farthest-first & 4/1, 0.62 & 3/4, 0.22 & 3/3, 0.84 & 5/3, 0.74 \\
Cov.-greedy & 1/2, 1.00 & 3/4, 0.58 & 2/4, 1.00 & 5/3, 0.95 \\
CLLM-style$^\dagger$ & 4/2, 1.00 & 2/6, 0.08 & 5/2, 0.38 & 3/5, 0.84 \\
Disc.-realism$^\dagger$ & 2/3, 0.81 & 4/4, 0.46 & 4/3, 0.94 & 3/5, 0.38 \\
Real only & 5/1, 0.44 & 3/5, 0.46 & 5/2, 0.08 & 5/3, 0.46 \\
CW-LR & 6/2, 0.31 & 3/5, 0.74 & -- & -- \\
ROS+$b$ & 5/1, 0.06 & 6/2, 0.38 & 5/1, 0.16 & 5/3, 0.38 \\
SMOTE(NC)+$b$ & 5/0, 0.06 & 3/5, 0.64 & 4/2, 0.56 & 6/2, 0.31 \\
SMOTE-ENN full & 4/2, 0.44 & 5/1, 0.09 & 3/3, 0.69 & 4/2, 0.31 \\
SMOTE-ENN +$b$ & 3/1, 0.62 & 3/3, 1.00 & 5/1, 0.16 & 5/1, 0.31 \\
SMOTE-Tomek full & 4/2, 0.44 & 4/2, 0.31 & 2/4, 0.31 & 3/3, 0.84 \\
SMOTE-Tomek +$b$ & 2/2, 1.00 & 2/4, 0.16 & 2/3, 0.81 & 4/2, 0.56 \\
B-SMOTE full & 4/2, 0.44 & 3/3, 0.56 & 3/3, 1.00 & 3/3, 0.69 \\
B-SMOTE +$b$ & 4/2, 0.56 & 2/4, 0.44 & 2/4, 1.00 & 3/3, 0.69 \\
ADASYN full & 5/1, 0.16 & 4/2, 0.22 & 2/4, 0.56 & 3/3, 0.69 \\
ADASYN +$b$ & 3/2, 0.81 & 3/2, 1.00 & 2/3, 1.00 & 3/2, 0.62 \\
\bottomrule
\end{tabular}}
\end{table}

Across all sixteen methods on those five datasets (Table~\ref{tab:addendum-ranks}), full-balancing SMOTE-Tomek and ADASYN have the best random-forest F1 ranks ($3.70$ and $4.50$, against $7.60$ for quality-coverage; Friedman $p=.031$), while no method dominates the other three classifier--metric combinations. No Wilcoxon comparison of quality-coverage with any method reaches $p<.05$ (smallest $.06$). Taken together, the addenda support the same reading as the main study on nine datasets: quality-coverage reliably improves the coverage and redundancy of the curated set, but it does not improve classifier utility over simple selection or published oversampling.
